% !PRINTED 60
\weberpage{60}
%<<CONT

\[
 \alpha_1+\alpha_2+\cdots+\alpha_m=\nu,
\]
the \gesperrt{derivatives of order $\nu$ of the function $\Phi$}.

They are also written, according to the notation usual in differential calculus, as
\begin{equation}
 D_{\alpha_1,\alpha_2,\ldots,\alpha_m}\Phi
 =
 {\partial^\nu\Phi\over
 \partial x_1^{\alpha_1}\partial x_2^{\alpha_2}\cdots\partial x_m^{\alpha_m}}.
\tag{5}
\end{equation}

The law of formation of the derivatives can be summarized in the following propositions, where for brevity we omit the indices on the symbol $D$.

\gesperrt{I. If $C$ is a constant, then}
\[
 D(C\Phi)=CD\Phi.
\]
\gesperrt{II. If $\Phi$ and $\Psi$ are any two functions, then}
\[
 D(\Phi+\Psi)=D\Phi+D\Psi.
\]
Both follow immediately from (4).

Thus we can easily form the derived functions in general if we know them for the special case in which $\Phi$ is a product of powers, that is, if we know
\[
 D_{\alpha_1,\alpha_2,\ldots,\alpha_m}(x_1^{\mu_1}x_2^{\mu_2}\cdots x_m^{\mu_m}),
\]
where the $\mu$ are arbitrary non-negative exponents.

But
\[
 (x_1+\xi_1)^{\mu_1}
 =\sum^{\alpha_1}{\Pi(\mu_1)\over\Pi(\alpha_1)\Pi(\mu_1-\alpha_1)}
 \xi_1^{\alpha_1}x_1^{\mu_1-\alpha_1},
\]
and therefore
\begin{equation}
\begin{aligned}
&(x_1+\xi_1)^{\mu_1}\cdots(x_m+\xi_m)^{\mu_m}\\
&=\sum^{\alpha_1\ldots\alpha_m}
 {\Pi(\mu_1)\cdots\Pi(\mu_m)x_1^{\mu_1-\alpha_1}\cdots x_m^{\mu_m-\alpha_m}
 \over
 \Pi(\mu_1-\alpha_1)\cdots\Pi(\mu_m-\alpha_m)\Pi(\alpha_1)\cdots\Pi(\alpha_m)}
 \xi_1^{\alpha_1}\cdots\xi_m^{\alpha_m}.
\end{aligned}
\tag{6}
\end{equation}
Comparison with (4) and (6) therefore gives
\begin{equation}
\begin{aligned}
&D_{\alpha_1,\ldots,\alpha_m}(x_1^{\mu_1}\cdots x_m^{\mu_m})\\
&\quad =
{\Pi(\mu_1)\cdots\Pi(\mu_m)\over
 \Pi(\mu_1-\alpha_1)\cdots\Pi(\mu_m-\alpha_m)}
 x_1^{\mu_1-\alpha_1}\cdots x_m^{\mu_m-\alpha_m},
\end{aligned}
\tag{7}
\end{equation}
as long as $\alpha_1\leq\mu_1,\ldots,\alpha_m\leq\mu_m$. On the other hand,
\begin{equation}
 D_{\alpha_1,\ldots,\alpha_m}(x_1^{\mu_1}\cdots x_m^{\mu_m})=0,
\tag{8}
\end{equation}
as soon as one of the indices $\alpha$ is greater than the corresponding exponent $\mu$.
